Let $A\subseteq \mathbb{N}$. Write $A(x):=\left|A\cap[1,x]\right|$ for $x\ge 1$ (with $[1,x]$ understood as $\{n\in\mathbb{N}:1\le n\le x\}$).
Assume:
\begin{align}
A(2x)-A(x) &\xrightarrow[x\to\infty]{} \infty,\label{eq:dyadic_growth}\\
\sum_{n\in A} \{\theta n\} &= \infty\quad\text{for every }\theta\in(0,1),\label{eq:theta_div}
\end{align}
where $\{y\}$ denotes the distance from $y$ to the nearest integer (so $\{y\}\in[0,1/2]$).

Conclusion to prove:
There exists $N_0$ such that every integer $N\ge N_0$ can be expressed as a sum of \emph{distinct} elements of $A$.
